\chapter{Dérivées, variations et extrémums}
\section*{Exercice 10.1 -- Variations d'un polynôme}
$f'(x)=3(x-1)(x+1)$ est positif sur $[-2,-1]$ et $[1,2]$, négatif sur $[-1,1]$. Les valeurs à comparer sont $f(-2)=-2$, $f(-1)=2$, $f(1)=-2$, $f(2)=2$. Le minimum global est $-2$ et le maximum $2$, chacun atteint en deux points.

\section*{Exercice 10.2 -- Dérivée nulle}
Non : pour $f(x)=x^3$, $f'(0)=0$ mais $f(x)<0$ à gauche et $f(x)>0$ à droite de $0$. Il n'y a aucun extremum en $0$. La condition est nécessaire à un extremum intérieur dérivable, mais pas suffisante.

\section*{Exercice 10.3 -- Parité}
$f(-x)=f(x)$, donc $f$ est paire. Sa dérivée $f'(x)=4x(x^2+1)$ satisfait $f'(-x)=-f'(x)$. Elle est négative pour $x<0$, nulle en $0$, positive pour $x>0$ ; $f$ décroît puis croît, avec minimum $0$.

\section*{Exercice 10.4 -- Hypothèse d'intervalle}
Non : $-1<1$ mais $f(-1)=-1<1=f(1)$, contraire à la décroissance globale. Elle est strictement décroissante sur chacun des intervalles $]-\infty,0[$ et $]0,+\infty[$. Le théorème reliant signe de la dérivée et monotonie s'applique sur un intervalle.

\section*{Exercice 10.5 -- Accroissements finis}
La dérivée de $\sqrt t$ vaut $1/(2\sqrt t)\le1/2$ sur $[1,+\infty[$ ; les accroissements finis donnent la borne. Algébriquement, $|\sqrt y-\sqrt x|=|y-x|/(\sqrt y+\sqrt x)\le|y-x|/2$. La seconde preuve est immédiate ici.

\section*{Exercice 10.6 -- Minimum en bord}
$f'(x)=2x+2>0$ sur $[0,3]$, donc $f$ croît et son minimum est $f(0)=0$. La dérivée ne s'annule pas dans l'intervalle. Un minimum situé à une extrémité n'est pas soumis à la condition de Fermat pour un point intérieur.
